\section{Certified selection by local response discovery}\label{sec:final_method}

\subsection{Design-selection problem and computational contract}

Let $m_0>0$ be the best observed objective, $\mathcal A$ a finite menu of first designs, and $\mathcal I_a$ the ordered continuation pool after first design $a$. The existing raw-response mean policy is
\begin{equation}
\pi_a(z)=\mathop{\rm argmin}^{\rm first}_{i\in\mathcal I_a}
                   (u_{ai}z+v_{ai}),\qquad z>0,
\label{eq:final_policy}
\end{equation}
where the smallest pool index resolves a tie. This policy is fixed during planning. The task is to select $a$ for the two-evaluation improvement
\begin{equation}
V_a=\mathbb E\left[m_0-\min\{m_0,Y_a,Y_{\pi_a(Y_a)}\}\right],
\qquad \max_{a\in\mathcal A}V_a.
\label{eq:final_objective}
\end{equation}
The optimization variable is the first design in the supplied menu, with the continuation policy held fixed.

The planning responses are positive. Write $\log Y_a=m_a+s_aX$, with $s_a>0$ and $X\sim N(0,1)$. When a continuation $i$ is encountered, a numerical query supplies enclosures for
\begin{equation}
\log Y_{ai}\mid X=x\ \sim\
N(\mu_{ai}+\gamma_{ai}x,\sigma_{ai}^2),\qquad \sigma_{ai}>0.
\label{eq:final_conditional_model}
\end{equation}
The raw-response model determines the continuation policy; the log-response model determines its planning reward.

Queries use the fixed training responses $y>0$, a fixed binary training matrix $A$, binary cross/target matrices $B,K$, and fixed binary normalization constants $c_0,s_0$. Their posterior arithmetic encloses
\begin{equation}
g=c_0\mathbf1+BA^{-1}(\log y-c_0\mathbf1),\qquad
S=s_0^2(K-BA^{-1}B^{\mathsf T}).
\label{eq:final_pair_posterior}
\end{equation}
For a queried first/second pair, $\mu_{ai}=g_2$, $\gamma_{ai}=S_{21}/s_a$ and $\sigma_{ai}^2=S_{22}-\gamma_{ai}^2$. The first marginal and pair use the same declared training inputs. Inverse residuals, positive-definiteness and positive conditional variance are checked before using a query. Parameters remain fixed during planning.

Let $\mathcal P_a$ denote the positive conditional-law completions agreeing with the queried marginals on certified regions. The certificate applies to every member of this class. A compatible joint GP has the same branch expectation and is included. Computation requires the queried marginals; the full unqueried joint covariance lies outside this contract.

\subsection{Certified local identity and regional value}

Let affine bands enclose the continuation means: $L_{ai}(z)\le u_{ai}z+v_{ai}\le U_{ai}(z)$ for all positive $z$. For reference action $j$, one scan intersects the linear constraints defining
\begin{equation}
D_{aj}=\{z>0: U_{aj}(z)<L_{ai}(z)\ (i<j),\quad
                  U_{aj}(z)\le L_{ai}(z)\ (i>j)\}.
\label{eq:final_stability}
\end{equation}
On this domain, the declared policy must select $j$. Earlier rivals require strict dominance to preserve the smallest-index tie rule on positive-measure intervals. Isolated endpoints have zero probability under the nondegenerate first-response model. Exact predicates support the linear comparisons \citep{ref41}. Each transformed endpoint is enclosed; using the upper bound for the lower endpoint and the lower bound for the upper endpoint gives a common inner standard-normal region valid for all enclosed $m_a,s_a$.

Define
\begin{equation}
h_a(x)=\min\{m_0,\exp(m_a+s_ax)\},\quad b_a(x)=m_0-h_a(x),
\quad w_{ai}(x)=\mathbb E[(h_a(x)-Y_{ai})_+\mid X=x].
\label{eq:final_remainder}
\end{equation}
The expected-improvement decomposition gives reward $b_a+w_{a\pi_a}$ with $0\le w_{ai}\le h_a$.

\paragraph{Proposition 1: regional value enclosure.}
Let $S_a$ be a union of certified response cells, with known policy identity on each cell, and $E_a=S_a^c$. Put $B_a=\mathbb E[b_a(X)]$, let $W_{a,S}$ be the integral of the appropriate known remainder over $S_a$, and put $H_{a,E}=\mathbb E[h_a(X)\mathbf1_{E_a}]$. Then
\begin{equation}
B_a+W_{a,S}\ \le V_a\le\ B_a+W_{a,S}+H_{a,E},
\qquad H_{a,E}\le m_0\mathbb P(E_a).
\label{eq:final_regional_enclosure}
\end{equation}
The unknown remainder lies between zero and $h_a$ pointwise, so integration proves the result for every admissible positive completion. The argument uses the pointwise range on $E_a$ and allows dependent candidate responses. If $E_a$ implies $Y_a\le\eta m_0$, then $H_{a,E}\le\eta m_0\mathbb P(E_a)$: an already good first response limits the remaining adaptation value.

For a standard-normal cell $C$, the implemented unknown-reward interval is $[B_C^-,m_0P_C^+]$, where $B_C=\int_C b_a\phi$ and $P_C=\int_C\phi$. A uniform-reward control uses $[0,m_0P_C^+]$. The capped-response moment is computed by splitting at $c=(\log m_0-m_a)/s_a$. Below $c$,
\begin{equation}
\int_r^t\exp(m_a+s_ax)\phi(x)\,dx
=\exp(m_a+s_a^2/2)\{\Phi(t-s_a)-\Phi(r-s_a)\};
\label{eq:final_moment}
\end{equation}
above it the capped moment is $m_0\{\Phi(t)-\Phi(r)\}$. An uncertain crossing strip receives a conservative reward hull. Acceptance uses the actual outward endpoints, including boundary and tail hulls.

\subsection{Conditional integration and numerical acceptance}

On an identified action/incumbent piece, write $h=\exp(\alpha+\beta x)$, with $\beta=s_a$ or $0$, $\mu=\mu_{ai}+\gamma_{ai}x$, $k=(\log h-\mu)/\sigma_{ai}$ and $M=\exp(\mu+\sigma_{ai}^2/2)$. The conditional reward is
\begin{equation}
q_i(x)=m_0-h\Phi(-k)-M\Phi(k-\sigma_{ai}).
\label{eq:final_q}
\end{equation}
The truncated-lognormal expectation is established \citep{ref45}. Differentiation using $h\phi(k)=M\phi(k-\sigma_{ai})$ gives the uniform smooth-piece bound
\begin{equation}
|q_i''|\le C_i=m_0\left\{\beta^2+\gamma_{ai}^2+
\frac{(\beta-\gamma_{ai})^2}{\sigma_{ai}\sqrt{2\pi}}\right\}.
\label{eq:final_curvature}
\end{equation}
For a cell $[r,t]$, normal-weighted linear interpolation therefore has absolute error at most $C_i(t-r)^2\{\Phi(t)-\Phi(r)\}/8$. Parameter intervals require a uniform bound on $C_i$ and a positive lower bound on $\sigma_{ai}$. The supplement gives the derivative and interpolation arguments.

The cells form a disjoint cover of $[-T,T]$, with $T=8$. Outside-window contribution is enclosed by $[0,2m_0\Phi(-T)]$, without renormalization. Directed logarithms, probabilities, interpolation, posterior residual bounds and final sums use MPFR arithmetic \citep{ref46}. Uncertain boundaries receive their full reward hull.

\paragraph{Proposition 2: accepted selection.}
For valid branch intervals, accept $\hat a$ only if
\begin{equation}
\max\{0,\max_{a\ne\hat a}\overline V_a-\underline V_{\hat a}\}
\le\delta,\qquad\delta=m_0/10^6.
\label{eq:final_acceptance}
\end{equation}
Every competitor is bounded above and the selected value below, so $\max_aV_a-V_{\hat a}\le\delta$ for every admissible positive completion. The competitor maximum excludes the selected action; a singleton menu is trivial. Failed arithmetic checks or an exhausted resource budget return unresolved.

\subsection{Algorithm and local error scope}

The algorithm returns an accepted first action and cell record, or unresolved:
\begin{enumerate}
\item Construct the raw affine bands and first-response enclosures from training inputs. Initialize the normal cell cover and split at the enclosed incumbent crossing. The input consists of training data, fixed parameters and the design menus.
\item Assemble current branch intervals. Choose the branch with largest lower bound and its competitor with largest upper bound. Apply Equation~\eqref{eq:final_acceptance} before requesting further work.
\item Select the widest value-uncertainty cell among these two branches. For an unknown cell, use a cached local domain if possible; otherwise select a reference at its center and perform the dominance scan in Equation~\eqref{eq:final_stability}. Split off the certified intersection, retaining unknown remainder cells.
\item Request a bivariate conditional posterior only for an encountered action, cache it, and integrate its certified cells. For an already identified cell, subdivide to reduce integration error. Return to the interval comparison until acceptance or the declared limit.
\end{enumerate}
The reader checks local dominance, tie order, cell coverage, fresh conditional integrals and sums. A separate audit reconstructs the tolerance from the input incumbent and checks regret. The producer and reader share arithmetic and curvature code.

Policy-composed reward need not be continuous. With $m_0=1$, raw means $z-1/2$ and $0$ switch the action at $z=1/2$. Independent lognormal planning responses with equal positive log-standard deviation and log-means $\log(1/4)$ and $\log(3/4)$ have strictly different expected capped minima at that point. The conditional reward consequently has different one-sided limits. Regional reward bounds remain valid across this switch.

For an idealized exact-parameter grid with maximum cell width $d$, let $J_a$ count policy switches and $C_a$ bound smooth-piece curvature. After a center probe per cell, unknown switch cells have total length at most $J_ad$. Their reward width, identified-cell interpolation width and normal tail give
\begin{equation}
\operatorname{width}([\underline V_a,\overline V_a])\le
\frac{m_0J_ad}{\sqrt{2\pi}}+\frac{C_ad^2}{4}+2m_0\Phi(-T).
\label{eq:final_local_error}
\end{equation}
This exact-parameter grid bound separates switch and smooth-integration errors. Fixed parameter bands and finite precision introduce additional floors, and the resource-limited schedule can stop before attaining a prescribed grid width. With $S$ local scans and $N$ eligible actions, dominance work is $O(SN)$ before exact-arithmetic bit costs, plus posterior and integration work.
